\documentclass{article}
\usepackage{graphicx} % Required for inserting images
\usepackage[spanish]{babel}
\usepackage{float}

\begin{document}

\begin{titlepage}
    \centering
    \vspace{5cm}
    {\Huge \textsc{\textbf{Práctica 1.: Grafana}}}\\
    \vspace{1cm}
    {\large \textsc{Evaluación de Sistemas Informáticos}}\\
    \vspace{1cm}
    {\large \textsc{3º Grado en Informática. Mención de TI}}\\
    \vspace{1cm}
    {\large \textsc{Grupo 6}}\\
    \vspace{1cm}
    {\large \textsc{Autores:}}\\
    {\large \textsc{Víctor Fernández Sánchez}}\\
    {\large \textsc{Gonzalo Sánchez Maroto}}\\
    {\large \textsc{Diego Sanz Santos}}\\
    {\large \textsc{Rubén Verduras Cavada}}\\
\end{titlepage}

\tableofcontents % Componer y mostar el índice
\listoffigures % Componer y mostrar el índice de tablas
\listoftables % Componer y mostrar el índice de figuras
\newpage

\section{¿Qué es Grafana?}

Grafana es un software de código abierto que permite consultar, visualizar, alertar y comprender datos en tiempo real sin importar dónde estén almacenados mostrados a través de paneles elegantes y flexibles.

Las principales características de Grafana son:

\begin{itemize}
    \item Permite crear dashboards altamente interactivos con gráficos y tablas que se pueden personalizar y explirar en tiempo real.
    \item Admite una amplia variedad de fuentes de datos como bases de datos SQL y NoSQL, sistemas de monitoreo como Prometheus y Graphite y sistemas en la nube como AWS CloudWatch y Google Stackdriver.
    \item Ofrece una amplia gama de opciones de visualización como gráficos de líneas, barras, dispersión, termómetros y más.
    \item Aporta la posibilidad de crear alertas basadas en umbrales y condiciones definidas por los usuario capaces de notificar a través de correo electrónico u otros canales de comunicación.
    \item Facilita la exploración y análisis de datos mediante funciones de búsqueda, filtrado y agregación, que permite a los usuarios investigar tendencias y patrones en los datos.
    \item Altamente seguro con sistemas basados en autenticación de usuarios y controles de acceso a datos basados en roles.
\end{itemize}

Según estas características algunos casos de uso son los siguientes:

\begin{enumerate}
    \item Monitoreo de infraestructura: supervisar el rendimiento de servidores, redes y bases de datos, identificando problemas y optimizando recursos.
    \item Visualización de métricas de aplicaciones: rastrea el rendimiento, identifica cuellos de botella y mejora la experiencia del usuario.
    \item Análisis de registros: correlacionar datos de registros con métricas para investigar problemas y tendencias en sistemas distribuidos.
    \item Observabilidad de microservicios: monitorizar la salud y rendimiento de estos para identificar problemas de comunicación entre servicios.
    \item Seguridad informática: visualizar datos de seguridad, como registros de acceso, eventos de autenticación y actividad sospechosa, para detectar y responder a las amenazas de manera más eficiente.
\end{enumerate}

\end{document}
